% ACM Artifact Evaluation appendix for HOLA.
\documentclass[manuscript,screen,nonacm]{acmart}

\usepackage{booktabs}
\usepackage{hyperref}
\usepackage{listings}

\settopmatter{printacmref=false,printfolios=true}
\setcopyright{none}
\title[HOLA Artifact]{HOLA Artifact Evaluation Appendix}
\author{Anonymous Authors}

\begin{document}
\maketitle
\appendix
\section{Artifact Appendix}

\subsection{Abstract}

This artifact contains the HOLA synthesizable reference model (SVM), the
NutShell RISC-V processor used as the design under test, a Linux boot image,
and an automated Artifact Evaluation workflow.  It reproduces the RCache
trade-off data in Figures 11--13, the phased Linux-boot plot in Figure 14, and
the RCore Chisel line counts in Table 4.  Every hardware experiment builds and
runs a Verilator simulator.  FPGA and Palladium platforms are not used.

\subsection{Artifact check-list (meta-information)}

{\small
\begin{itemize}
  \item {\bf Algorithm:} HOLA hostless co-simulation, RCore, and RCache.
  \item {\bf Program:} SVM and NutShell, written in Chisel/Scala, SystemVerilog,
        C, C++, Python, and shell.
  \item {\bf Compilation:} Mill, Chisel, LLVM/Clang or GCC, and Verilator.
  \item {\bf Binary:} The repository includes the NutShell Linux boot image;
        all simulator binaries are built by the workflow.
  \item {\bf Data set:} NutShell \texttt{ready-to-run/linux.bin}.
  \item {\bf Run-time environment:} OpenXiangShan x86-64 Ubuntu 24.04 Docker
        container.
  \item {\bf Hardware:} A 16-core x86-64 host with at least 32 GiB RAM is
        recommended; no accelerator or programmable logic is required.
  \item {\bf Execution:} Serialized emulator builds and parallel Verilator runs.
  \item {\bf Metrics:} Evict miss rate, instruction miss rate, refill bandwidth
        overhead, IPC, load/store IPC, and physical Chisel lines of code.
  \item {\bf Output:} Four text tables, one PDF plot, one LoC table, and raw logs.
  \item {\bf Experiments:} MICRO 2026 paper Figures 11--14 and Table 4.
  \item {\bf Disk space required:} Approximately 20 GiB including the Docker
        image, dependency caches, build products, and raw logs.
  \item {\bf Preparation time:} Approximately 20--40 minutes, depending on
        network and package-cache speed.
  \item {\bf Experiment time:} Approximately 2--4 hours on a 16-core host.
  \item {\bf Publicly available:} The source archive is supplied with the
        artifact submission; the two pinned source components are public GitHub
        repositories.
  \item {\bf Code licenses:} SVM and NutShell license files are included in their
        respective submodules.
  \item {\bf Workflow automation:} Fail-fast Bash scripts plus Python reporting.
\end{itemize}
}

\subsection{Description}

\subsubsection{How to access}

Obtain the artifact as a Git checkout and initialize its submodules
recursively.  The evaluated revisions are
NutShell commit
\texttt{c1394855f28b301f5b2cb835ffa642e97bec43c4} from branch \texttt{svm}
and SVM commit
\texttt{0c5299a55f9a9c57acc81be36467b18eb7faa96f} from branch \texttt{main}.
Clone the repository recursively when an artifact-service URL is available, or
enter an existing Git checkout and initialize all nested sources:

\begin{lstlisting}[language=bash]
git clone --recursive ARTIFACT_REPOSITORY_URL micro2026ae
cd micro2026ae
git submodule update --init --recursive
\end{lstlisting}

\subsubsection{Hardware dependencies}

The workflow requires an x86-64 machine capable of running Docker.  We
recommend at least 16 logical CPU cores, 32 GiB of memory, and 20 GiB of free
disk space.  The default run concurrency is four Verilator processes.  The
experiments do not use an FPGA, Palladium, VCS, or any commercial EDA tool.

\subsubsection{Software dependencies}

The only host dependency is Docker.  The OpenXiangShan xs-env image provides
the compiler toolchain, OpenJDK, Mill, and Verilator 5.048 or newer.  Inside the
container, \texttt{setup/setup.sh} verifies those tools and installs only
missing supporting packages such as Matplotlib.  It invokes the upstream
\texttt{make init} targets in NutShell and SVM; SVM initializes its instruction
semantics from pinned Spike commit
\texttt{055624200a34bcc8d4c7acde\allowbreak040466f7d2001637}.

\subsubsection{Data sets}

All simulations boot \texttt{NutShell/ready-to-run/linux.bin}.  The workload
boots Linux through the \texttt{/init} process, prints
\texttt{Hello, RISC-V World!}, and exits through the simulator good-trap path.
The image is included in the NutShell submodule; no external data download is
needed after repository setup.

\subsection{Installation}

Pull the evaluated OpenXiangShan Ubuntu 24.04 image:

\begin{lstlisting}[language=bash]
docker pull ghcr.io/openxiangshan/xs-env:ubuntu-24.04
\end{lstlisting}

From the directory containing \texttt{micro2026ae/}, run the complete workflow
in a fresh container.  The repository is mounted so that results remain on the
host after the container exits.

\begin{lstlisting}[language=bash]
docker run --rm --name micro2026ae \
  -v "$PWD/micro2026ae:/workspace/micro2026ae" \
  -w /workspace/micro2026ae \
  ghcr.io/openxiangshan/xs-env:ubuntu-24.04 \
  bash -lc './run-all.sh'
\end{lstlisting}

The container needs network access during setup for any missing Ubuntu
packages, Spike, and uncached Scala dependencies.  No network access is
required after
\texttt{setup/setup.sh} has completed successfully.

\subsection{Experiment workflow}

\texttt{run-all.sh} invokes the following four fail-fast stages:

\begin{enumerate}
  \item \texttt{setup/setup.sh} validates dependencies, initializes both
        upstream projects through their \texttt{make init} targets, and checks
        all fixed revisions.
  \item \texttt{build/build.sh} elaborates NutShell once, then serially builds
        each unique RCache configuration.  Every temporary source tree,
        generated RTL, compiler output, emulator, and log lives below a
        configuration-specific \texttt{build/build-*} directory.  Successful
        configurations are fingerprinted and can be resumed safely.
  \item \texttt{run/run.sh} boots Linux on every emulator.  Independent runs
        execute in parallel and write raw logs below \texttt{results/raw}; the
        requested \texttt{results/figure11} through
        \texttt{results/figure14} directories contain links to the relevant
        raw points.  A run is accepted only after a good trap, except for the
        documented 256 KiB/4-way capacity abort in Figure 11.
  \item \texttt{report/report.sh} derives all metrics from cumulative counter logs,
        validates the Figure 14 time series, invokes
        \texttt{SVM/scripts/plot.py}, and counts Table 4 directly from the Scala
        source.
\end{enumerate}

The stages may also be run separately:

\begin{lstlisting}[language=bash]
bash setup/setup.sh
bash build/build.sh
bash run/run.sh
bash report/report.sh
\end{lstlisting}

\subsection{Evaluation and expected results}

A successful complete run creates the following files:

\begin{lstlisting}
evaluation/figure11.txt
evaluation/figure12.txt
evaluation/figure13.txt
evaluation/figure14.txt
evaluation/figure14.pdf
evaluation/table4.txt
\end{lstlisting}

Each text report contains both measured values and the paper values, followed
by the maximum absolute difference in percentage points.  The expected paper
data are shown below.

\paragraph{Figure 11.}
Evict miss rate is
$100\times\mathtt{cache\_evict\_miss}/\mathtt{cache\_evict}$.

\begin{center}
\begin{tabular}{lrrr}
\toprule
Cache & 4-way & 8-way & 16-way \\
\midrule
256 KiB & X & 2.76\% & 3.42\% \\
512 KiB & 0.60\% & 0.41\% & 0.36\% \\
1 MiB   & 0.06\% & 0.01\% & 0.00\% \\
2 MiB   & 0.02\% & 0.00\% & 0.00\% \\
\bottomrule
\end{tabular}
\end{center}

The X is an expected abort because every block becomes dirty and
non-replaceable.  The supplementary PLRU point reported in the paper text is
0.53\% for the 256 KiB/8-way configuration.

\paragraph{Figure 12.}
Instruction miss rate is
$100\times\mathtt{core\_out\_miss}/\mathtt{core\_out}$.

\begin{center}
\begin{tabular}{lrrrr}
\toprule
Configuration & 256 KiB & 512 KiB & 1 MiB & 2 MiB \\
\midrule
2-bank, 2-port & 18.32\% & 2.08\% & 0.20\% & 0.19\% \\
2-bank, 3-port & 18.18\% & 1.89\% & 0.02\% & 0.00\% \\
4-bank, 2-port & 18.27\% & 2.00\% & 0.12\% & 0.11\% \\
4-bank, 3-port & 18.18\% & 1.89\% & 0.01\% & 0.00\% \\
\bottomrule
\end{tabular}
\end{center}

\paragraph{Figure 13.}
The disabled and enabled rows use the Figure 12 miss metric.  Bandwidth
overhead is
\[
100\times
\frac{\mathtt{cache\_read\_cache\_miss}}
     {\mathtt{cache\_refill}+\mathtt{cache\_evict}}.
\]

\begin{center}
\begin{tabular}{lrrrr}
\toprule
Series & 256 KiB & 512 KiB & 1 MiB & 2 MiB \\
\midrule
Disabled & 18.32\% & 2.08\% & 0.20\% & 0.19\% \\
Enabled  & 0.70\% & 0.26\% & 0.19\% & 0.19\% \\
Overhead & 2.70\% & 0.41\% & 0.00\% & 0.00\% \\
\bottomrule
\end{tabular}
\end{center}

\paragraph{Figure 14.}
\texttt{evaluation/figure14.pdf} contains exactly the paper's four displayed
series: 2-port miss rate, 3-port miss rate, IPC, and load/store IPC.  The plot
should cover approximately $300\times2^{16}$ reference cycles and show the
large early 2-port spikes removed by the third port.  The plotting command does
not add the other diagnostic performance-counter series available in
\texttt{plot.py}.

Both Figure 14 runs use a 512 KiB, 8-way cache with refill-on-read-miss.  The
2-port point uses two banks; the 3-port point uses the paper's four-bank,
three-port design.

The evaluated build keeps the RTL's 1,024-cycle counter snapshots but disables
the bootrom's final UART counter dump.  This is the no-performance-counter-output
variant used for Figure 14; it avoids adding a long diagnostic dump after the
Linux workload while retaining the samples required by \texttt{plot.py}.

\paragraph{Table 4.}
The expected module counts are 253, 61, 52, 126, 87, 154, 49, 25, and 129
lines for MMU, FETCH, AGU, LSU, ARITH, PRIV, TRAP, COMMIT, and RCore,
respectively, totaling 936 physical Chisel source lines.  The report excludes
the later Figure 14-only instrumentation and fails if any source boundary has
changed.

\subsection{Experiment customization}

Set \texttt{AE\_RUN\_JOBS=N} to change simulator concurrency.  Set
\texttt{AE\_RUN\_TIMEOUT} to change the per-run timeout in seconds (default:
3,600).  A
space-separated subset of configuration identifiers may be supplied through
\texttt{AE\_CONFIGS}; identifiers are listed in
\texttt{experiments/configs.tsv}.  For example:

\begin{lstlisting}[language=bash]
AE_CONFIGS="c512k-w8-b2-p2-r1 c512k-w8-b4-p3-r1" \
  bash build/build.sh
AE_CONFIGS="c512k-w8-b2-p2-r1 c512k-w8-b4-p3-r1" \
  AE_RUN_JOBS=2 bash run/run.sh
bash report/report-figure14.sh
\end{lstlisting}

Set \texttt{FORCE\_REBUILD=1} or \texttt{FORCE\_RERUN=1} to ignore successful
fingerprints.  Set \texttt{KEEP\_BUILD\_WORK=1} to retain complete temporary
source and compiler trees for debugging.  \texttt{FIGURE14\_XMAX} changes the
plot's maximum x-axis value (default 305).

\subsection{Notes}

All scripts use \texttt{set -Eeuo pipefail}; an unexpected dependency,
compilation, simulation, or reporting failure stops the complete workflow.
Builds are intentionally serialized because both upstream projects normally
write fixed \texttt{build/} paths.  This artifact avoids interference by using
configuration-specific source copies.  Runs use separate emulator binaries and
result directories and are therefore parallel-safe.

If setup cannot download dependencies, first verify that the Docker container
has outbound HTTPS and DNS access.  If a simulator fails, inspect
\path{results/raw/CONFIG/stdout.log} and \path{stderr.log}.  If a build
fails, inspect \path{build/build-CONFIG/build.log}; rerun that configuration
with \texttt{KEEP\_BUILD\_WORK=1} to retain its generated source tree.

\subsection{Methodology}

Submission, reviewing, and badging follow the ACM policy:
\begin{itemize}
  \item \url{https://www.acm.org/publications/policies/artifact-review-and-badging-current}
  \item \url{https://cTuning.org/ae}
\end{itemize}

\end{document}
